% SPDX-License-Identifier: CC-BY-SA-4.0
% Copyright (c) 2026 Luis Alonso
%
% This file is part of luis_alonso/Notas-Japones.
%
% Licensed under Creative Commons Attribution-ShareAlike 4.0 International (CC BY-SA 4.0).
% You may share and adapt this material for any purpose,
% as long as you provide attribution and distribute adaptations under the same license.
%
% License: https://creativecommons.org/licenses/by-sa/4.0/
% Source:  https://git.lalonso.com/luis_alonso/Notas-Japones
%
% Note: Third-party content (if any) is excluded from this license and is marked where used.

\documentclass[a4paper,lualatex,ja=standard,base=10pt]{bxjsarticle}
\usepackage{setup}
\usepackage{ulem}

\newcommand{\rmv}[1]{\textcolor{red}{\sout{#1}}}

\title{\ruby{第|０５|課}{だい||か} LV4}
\author{アロンゾ　ルイス}

\begin{document}
\setcounter{page}{1}

\maketitle{}
\thispagestyle{fancy}

\section{Dar mas de un adjetivo a un sujeto (2)}

Para dar multiples cualidades a un sujeto (adjetivos) se usa くて o で con la siguiente formula:

\begin{itemize}
  \item イAdj pasa a -くて
  \item ナAdj pasa a -で
\end{itemize}

Se usa para concatenar adjetivos con el mismo sentido (ambos positivos o ambos negativos).

\section{Enlistar cosas que hay en lugres}

Para enlistar cosas pero no una lista exhaustiva se usa la particula や en lugar del と.

El と se usa para listas (completas) exhaustivas.

\section{Contraste con けど}

Para hacer contraste entre los adjetivos enlistados se usa la particula けど.

\begin{itemize}
  \item Para イAdj se agrega けど al primer adjetivo y el segundo queda intacto.
  \item Para ナAdj se agrega だけど al primer adjetivo y el segundo queda intacto.
\end{itemize}

Para esta forma el adjetivo con mas peso es el ultimo.

\section{Vocabulario}

\begin{itemize}
  \item すすめ Recomendación
  \item ばしょ Lugar (para recomendar)
  \item \ruby{所}{ところ} Lugar (para recomendar)
  \item \ruby{市場}{いちば} Mercado
  \item こうきょ Palacio imperial
  \item このあたり En esta area
  \item どのあたりですか。 Por donde esta?
  \item ゆうえんち Parque de diversiones
  \item 雨だけど、行きます Aunque este lloviedo ve.
\end{itemize}

\end{document}

% a
% b
% c
